\subsection{Planteo y Despeje de Ecuaciones Lineales}

\graybold{Preguntas guía:}

\begin{itemize}
\item ¿El enunciado describe una relación entre cantidades que crecen a la misma tasa (por ejemplo, "años transcurridos desde X")?
\item ¿Se pide encontrar un punto en el tiempo (o un valor) donde se cumple una proporción exacta entre esas cantidades?
\item ¿La relación se puede plantear como una única ecuación lineal con una incógnita?
\end{itemize}

\graybold{¿Para qué sirve?} Reconocer cuándo un problema, a pesar de su narrativa, se reduce a plantear una ecuación lineal simple y despejar — sin necesidad de ningún algoritmo, búsqueda o simulación. Es una habilidad tan importante como reconocer cuándo SÍ hace falta algo más complejo.

\graybold{Complejidad:} \bigO{1}.

\graybold{Fundamento teórico:} Si el maestro empezó en el año $W$ y el aprendiz en el año $K$, sus años de experiencia en un año $Y$ cualquiera son $Y-W$ y $Y-K$ respectivamente. Pedir que la experiencia del maestro sea EXACTAMENTE el doble se traduce en $Y - W = 2(Y-K)$, que al despejar da $Y = 2K - W$ — una fórmula cerrada, sin necesidad de iterar ni buscar.

\graybold{Ejercicio aplicado}

Dados los años en que el maestro Wei ($W$) y su aprendiz Kai ($K$) comenzaron a programar, hallar el año exacto en que la experiencia del maestro será el doble de la del aprendiz.

\graybold{I/O:} Una sola línea con dos enteros → lectura directa, sin ninguna consideración especial de rendimiento.

\graybold{Código solución}

\begin{lstlisting}[language=Python]
import sys

def solve():
    w, k = map(int, sys.stdin.buffer.read().split())
    # se busca Y tal que (Y - w) = 2*(Y - k)  =>  Y = 2k - w
    print(2 * k - w)

solve()
\end{lstlisting}